\section{O`zgarmas tok}



\begin{enumerate}

\item 800 V potensiallar farqigacha zaryadlangan havo kondensatori o‘lchamlari jihatidan birday bo‘lgan, dielektrik to‘ldirilgan zaryadlanmagan kondensator bilan parallel ulangan. Agar kondensatorlar ulangandan keyin potensiallar farqi 100 V ga teng bo‘lgan bo‘lsa, dielektrikning singdiruvchanligini aniqlang.
\begin{enumerate}[label=\Alph*)]
\item 7
\item 2
\item 16
\item 3
\end{enumerate}

\item Yassi havo kondensatori 50 V potensiallar farqigacha zaryadlangan va manbadan uzib qo‘yilgan. So‘ngra kondensator qoplamalari orasiga, ularga parallel qilib, 1 mm qalinlikdagi metall listi kiritiladi. Agar qoplamalar orasidagi masofa 5 mm bo‘lsa, kondensator qoplamalari orasidagi potensiallar farqini aniqlang (V).
\begin{enumerate}[label=\Alph*)]
\item 30
\item 15
\item 80
\item 40
\end{enumerate}

\item Ikkita bir xil kondensator parallel ulangan va 6 V potensiallar farqigacha zaryadlangan. Agar kondensatorlar manbadan uzilgandan so‘ng ulardan birining plastinkalari orasidagi masofa ikki baravar kamaytirilgan bo‘lsa, kondensatorlar plastinkalari orasidagi potensiallar farqini aniqlang (V).
\begin{enumerate}[label=\Alph*)]
\item 6
\item 2
\item 4
\item 12
\end{enumerate}

\item O‘zgarmas tok manbaiga ulangan 12 $\mu F$ sig‘imli kondensatorga parallel ulangan voltmetr 3 V kuchlanishni ko‘rsatdi. Agar bu kondensatorga 8 $\mu F$ sig‘imli ikkinchi kondensator parallel ulansa, voltmetr qanday kuchlanishni (V) ko‘rsatadi?
\begin{enumerate}[label=\Alph*)]
\item 1,5
\item 3
\item 1
\item 2
\end{enumerate}

\item Zaryadli zarralarning tartibli harakatiga nima deyiladi?
\begin{enumerate}[label=\Alph*)]
\item tok kuchi
\item tok zichligi
\item elektr toki
\item elektr qarshilik
\end{enumerate}

\item Elektr tokining asosiy belgisi nima?
\begin{enumerate}[label=\Alph*)]
\item Uning atrofida magnit maydonning mavjudligidir.
\item Uning atrofida gravitatsion maydonning mavjudligidir.
\item Uning atrofida potensial maydonning mavjudligidir.
\item Uning atrofida uyurmaviy maydonning mavjudligidir.
\end{enumerate}

\item To‘g‘ri tasdiqni ko‘rsating.
\begin{enumerate}[label=\Alph*)]
\item Tok kuchi vektor kattalikdir.
\item Tok kuchi skalyar va musbat kattalik.
\item Tok kuchining ishorasi zaryad tashuvchi zarralar ishorasi bilan aniqlanadi.
\item Tok kuchi skalyar, ammo ishorali kattalik.
\end{enumerate}

\item Tok zichligi deb nimaga aytiladi?
\begin{enumerate}[label=\Alph*)]
\item o‘tkazgich sirtidan vaqt birligi ichida o‘tuvchi...
\item o‘tkazgich bo‘ylab vaqt birligida 1 m ga siljigan...
\item o‘tkazgichning ko‘ndalang kesimidan vaqt birligi ichida o‘tuvchi...
\item o‘tkazgich ko‘ndalang kesimi yuza birligidan vaqt birligida o‘tuvchi zaryad miqdoriga teng vektor kattalik.
\end{enumerate}

\item O‘tkazgichdan 10 A tok oqmoqda. Xuddi shu o‘tkazgichning ko‘ndalang kesimidan 1 soat davomida oqib o‘tadigan elektronlarning umumiy massasini toping (mg). $e = 1,6 \cdot 10^{-19} \, C, m_e = 9,1 \cdot 10^{-31} \, kg$.
\begin{enumerate}[label=\Alph*)]
\item 0,4
\item 0,8
\item 0,2
\item 0,12
\end{enumerate}

\item Tok kuchi 32 $\mu A$ bo‘lganda 1 ns vaqt ichida o‘tkazgichning ko‘ndalang kesimidan qancha elektron o‘tadi?
\begin{enumerate}[label=\Alph*)]
\item $1 \cdot 10^5$
\item $2 \cdot 10^3$
\item $4 \cdot 10^5$
\item $2 \cdot 10^5$
\end{enumerate}

\item Tok kuchi 16 mA bo‘lganda o‘tkazgich ko‘ndalang kesimidan 1 ms da o‘tadigan elektronlar sonini aniqlang. $e = 1,6 \cdot 10^{-19} \, C$.
\begin{enumerate}[label=\Alph*)]
\item $6 \cdot 10^{14}$
\item $5,6 \cdot 10^{14}$
\item $5 \cdot 10^{14}$
\item $10^{14}$
\end{enumerate}

\item O‘tkazgichdan 24 soat davomida 1 A tok o‘tgan. Bu davrda o‘tkazgich kesimi orqali o‘tgan elektronlar massasi (mg) topilsin. $e/m_e = 1,76 \cdot 10^{11} \, C/kg$.
\begin{enumerate}[label=\Alph*)]
\item 0,5
\item 2
\item 5
\item 0,3
\end{enumerate}

\item Agar o‘tkazgichdan o‘tayotgan tok kuchi 10 s da 0 dan 5 A gacha chiziqli ortgan bo‘lsa, shu vaqt ichida o‘tkazgich ko‘ndalang kesimidan qancha zaryad oqib o‘tgan (C)?
\begin{enumerate}[label=\Alph*)]
\item 0,5
\item 2,5
\item 2
\item 25
\end{enumerate}

\item Rux sulfat solingan elektrolitik vannada tok kuchi chiziqli $I = (2 + 0,02t) \, A$ qonun bo‘yicha o‘zgaradi. Elektrolit orqali 5 min ichida o‘tgan elektr zaryadni aniqlang (C)?
\begin{enumerate}[label=\Alph*)]
\item 1400
\item 1500
\item 1600
\item 1800
\end{enumerate}

\item Zanjirdagi tok kuchi $I = 10^{-6}t$ qonun bo‘yicha ortib bormoqda. Zanjirdan 1,5 minut davomida qancha elektron o‘tgan?
\begin{enumerate}[label=\Alph*)]
\item $2,5 \cdot 10^{16}$
\item $2,6 \cdot 10^{14}$
\item $3,9 \cdot 10^{17}$
\item $2,53 \cdot 10^{16}$
\end{enumerate}

\item Zanjirning bir qismidagi tok kuchi 6 s ichida 0 dan 1,5 A gacha bir tekis o‘sdi. Tok kuchining vaqtga bog‘lanish grafigini ko‘rsating (Source [3] dagi grafiklar asosida).
\begin{enumerate}[label=\Alph*)]
\item 1-grafik (3 s da 1,5 A)
\item 2-grafik (9 s da 1,5 A)
\item 3-grafik (6 s da 3 A)
\item 4-grafik (12 s da 3 A)
\end{enumerate}

\item Sig‘imi 200 $\mu F$ bo‘lgan kondensator 5 s vaqt ichida 5000 V kuchlanishgacha zaryadlanadi. Zaryad tok kuchining o‘rtacha qiymatini (A) toping.
\begin{enumerate}[label=\Alph*)]
\item 0,02
\item 0,05
\item 2
\item 0,2
\end{enumerate}

\item 8 s vaqt ichida o‘tkazgichning 5 $mm^2$ ko‘ndalang kesimidan 20 C zaryad o‘tgan bo‘lsa, tok zichligi qanday ($A/m^2$)?
\begin{enumerate}[label=\Alph*)]
\item $5 \cdot 10^5$
\item $4 \cdot 10^4$
\item $4 \cdot 10^5$
\item $4 \cdot 10^6$
\end{enumerate}

\item Ko‘ndalang kesimi 20 $mm^2$ bo‘lgan mis sim orqali 3 A tok o‘tayotgan bo‘lsa, undagi o‘tkazuvchanlik elektronlarining tezligi ($m/s$) topilsin. Birlik hajmdagi erkin elektronlarning soni $9 \cdot 10^{28} \, m^{-3}$ ga teng.
\begin{enumerate}[label=\Alph*)]
\item $3 \cdot 10^{-5}$
\item $10^{-5}$
\item 3
\item $2 \cdot 10^{-5}$
\end{enumerate}

\item Ko‘ndalang kesim yuzi 0,5 $mm^2$ bo‘lgan metall o‘tkazgichdan 3 A tok o‘tadi. Agar metalldagi erkin elektronlarning konsentratsiyasi $4 \cdot 10^{28} \, m^{-3}$ bo‘lsa, elektronlar dreyfining o‘rtacha tezligi topilsin ($mkm/s$).
\begin{enumerate}[label=\Alph*)]
\item 1,2
\item 4,7
\item 2,8
\item 9,4
\end{enumerate}


\end{enumerate}



\begin{table}[!ht]
\centering
\begin{tabular}{|c|c|c|c|c|c|c|c|c|c|c|c|}
\hline
№ & J. & № & J. & № & J. & № & J. & № & J.  \\ \hline
1 & A & 5 & C & 9 & C & 13 & D & 17 & D  \\ \hline
2 & D & 6 & A & 10 & D & 14 & B & 18 & A  \\ \hline
3 & C & 7 & D & 11 & D & 15 & D & 19 & B  \\ \hline
4 & B & 8 & D & 12 & A & 16 & D & 20 & D  \\ \hline
\end{tabular}
\end{table}